\begin{exercise}{16}
Show that the following simplicial sets are not nerves of categories:
\begin{enumerate}[label=(\arabic*),leftmargin=*,itemsep=0pt,parsep=0pt,
                  topsep=.5\baselineskip]
 \item $\partial\Delta^n$ for $n\ge2$
 \item $\Lambda^n_j$ for $n\ge2$ and $0\le j\le n$, except
 $(n,j)=(2,0),(2,2)$\footnote{The original gives the range
 $n\ge2$ and $0\le j\le n$ without exceptions. The corrected range
 excludes the two outer $2$-horns: $\Lambda^2_0$ and $\Lambda^2_2$
 are the nerves of the categories $1\leftarrow0\to2$ and
 $0\to2\leftarrow1$, respectively. Each category has only two
 nonidentity arrows and no composable pair of nonidentity arrows,
 so its only nondegenerate simplices are the vertices and edges of
 the stated horn.}
 \item $I^n$ for $n\ge2$
\end{enumerate}

\begin{proof}
By Definition~1.1.43, every string of $n$ composable arrows in a nerve
extends to an $n$-simplex. Suppose $Z\subseteq\Delta^n$ contains
$I^n$ and is a nerve. Its successive edges
$0\to1\to\cdots\to n$ must then extend to an $n$-simplex of $Z$
with vertex list $(0,1,\ldots,n)$. The only simplex of $\Delta^n$
with this list is $\mathrm{id}_{[n]}$. Thus any such $Z$ omitting
$\mathrm{id}_{[n]}$ cannot be a nerve.

For $n\ge2$, both $\partial\Delta^n$ and $I^n$ contain the spine
and omit $\mathrm{id}_{[n]}$: every spine simplex has at most two
vertices, hence is not surjective onto $[n]$.
Exercise~\ref{ex:13} gives the same obstruction for all
$\Lambda^n_j$ with $n\ge3$ and for $\Lambda^2_1$.
This proves all the corrected assertions.
\end{proof}
\end{exercise}
